% ===========================================================================
\section{Related work}
\label{sec:related}

The QEC mechanisms this paper touches are established, and we state the overlaps specifically.

\paragraph{Noise calibration from detector statistics.} Calibrating noise models from pairwise detector correlations, through the $p_{ij}$ estimator and the edge-weight formula $w_{ij} = \log((1-p_{ij})/p_{ij})$, is standard experimental practice~\cite{fowler2014scalable,chen2021exponential}, deployed at scale in Google's surface-code experiments~\cite{google2023suppressing,acharya2025belowthreshold}. Detector statistics have likewise been used to detect and remove correlated and leakage errors~\cite{mcewen2021removing}, and hidden-Markov inference on repeated syndrome records has been used to detect leakage~\cite{bultink2020leakage,varbanov2020hmm}. E9's recalibration is simpler than any of these: it matches the mean detector-flip rate to a simulated curve for the one channel known to drift and rebuilds the decoder from the chosen candidate's detector error model. Estimating the noise model from syndrome statistics alone, without disturbing the logical information, has its own literature: in-situ estimation proposals~\cite{fujiwara2014instantaneous,combes2014insitu}, identifiability results~\cite{wagner2021optimal,wagner2022paulichannels}, and learning of detector models and correlated noise from experimental syndrome data~\cite{chen2022calibrated,harper2023correlated,gicev2024syndrome,remm2025syndromecorrelations,blumekohout2025dem}.

\paragraph{Decoders under drift.} Reweighting a matching decoder's graph in response to drifting noise is the direct subject of prior work~\cite{spitz2018adaptive,huo2017realtime,dgr2023reweighting}, as is calibrating decoder priors to correlated or device-specific noise~\cite{nickerson2019correlated,sivak2024priors}. Among the closest neighbours of Section~\ref{sec:arc}'s drift witness that our bounded literature search found are the time-varying syndrome estimators of~\cite{kobori2025bayesian,bhardwaj2025drifting} and the drift-triggered recalibration of~\cite{poster2026reloqate}. In-situ recalibration of physical controls during repeated error detection was demonstrated in~\cite{kelly2016insitucalibration}. Drift itself is a characterized phenomenon~\cite{proctor2020drift}, and sequential change detection by CUSUM is classical~\cite{page1954continuous}.

\paragraph{Budgeted syndrome measurement.} Trading measurement budget against decoding quality by measuring fewer, or adaptively chosen, checks is also established~\cite{berthusen2024partialsyndrome,tansuwannont2023adaptivesyndrome,berthusen2025adaptivesyndrome}, though not with a coverage-residual objective.

\paragraph{Tools.} Our simulation and decoding infrastructure is Stim~\cite{gidney2021stim} and PyMatching~\cite{higgott2022pymatching,higgott2025sparseblossom}, used as documented libraries.

\paragraph{What this paper adds.} Relative to all of this, we claim no novelty of mechanism. What we add is orthogonal. First, a certificate framing in which coverage, existence, memory and price are typed audit objects derived from one framework. Second, an exact-computation layer, notably the Walsh--Hadamard inversion of the syndrome distribution behind Section~\ref{sec:diagnosis}, used after E2 had failed to \emph{explain} why a detector-statistic method failed, and then to predict in advance where its corrected form would and would not win. Third, a pre-registration discipline under which the framework's predictions, including predictions of its own failure, were frozen and scored.

\paragraph{The SBT side.} This paper takes the SBT papers~\cite{sbt-f1,sbt-f3,sbt-xi,sbt-hol,sbt-pt,sbt-ng,sbt-cast,sbt-xor,sbt-spend,sbt-qm,sbt-hid,sbt-nk} as its specification and tests them; it does not extend the theory. Appendix~\ref{app:formulas} cites the instantiated formulas object by object, and Section~\ref{sec:sbt} notes, primitive by primitive, the established external concepts (lumpability, causal states, projection-operator memory) that the instantiated objects correspond to.
